\scanpage{25}

On voit donc, grâce à la proposition 3, la validité des deux premiers énoncés contenus dans

\uline{Prop} 7\quad Soit $X$ un espace loc. noeth. spectral, $\mathcal F$ une catégorie fibrée sur la catégorie des ouverts de $X$, et \deleted{soit} $Y$ une partie fermée de $X$. \deleted{\unsure{On}} Conditions équivalentes

(i) $\mathcal F$ satisfait à $P_1$ ($P_2,P_3$) rel. à $\unsure{x|Y}$

(ii) Pour tt $x\in Y$, $\mathcal F$ satisfait à $P_{\deleted{\unsure{3}}1}$ ($P_2,P_3$) en $x$.

Reste le cas $P_3$, \deleted{\unsure{donc à prouver}} quand $P_1,P_2$ sont déjà vérifiés. On va prouver le

\margin{$\unsure{\Rightarrow|3}$}

\deleted{$(\mathrm{i})\Rightarrow\unsure{(\mathrm{ii})|\mathrm{ii}|(\mathrm{i})}$}

\uline{Corollaire} \uuline{8}\quad Si $\mathcal F$ satisfait $P_1$ ($P_2,P_3$) rel à $Y$, alors pour \deleted{tt ouvert $X'\subset X$, et} tte partie fermée $Y'\subset \deleted{\unsure{X'|Y}}\unsure{Y|X'}$ \deleted{telle que $Y'\subset Y\cap X'$}, $\mathcal F\deleted{\mid X'}$ satisfait $P_1$ ($P_2,P_3$) rel. à $Y'$.

Déjà fait pour $P_1,P_2$ voyons $P_3$. Voyons le diagramme

% archive sheet number: 24
